\section{Solutions for Chapter~\ref{ch:sewing}}

\paragraph{Exercise~\ref{ex:sewing-Ising}.}
\phantomsection
\label{sol:ex:sewing-Ising}
For $0<z<1$, put $u=\sqrt{1+\sqrt{1-z}}$ and
$v=\sqrt{1-\sqrt{1-z}}$. Then
$u^2+v^2=2$ and $uv=\sqrt z$.
Consequently
$(u\pm v)/\sqrt2=\sqrt{1\pm\sqrt z}$.
The common factor $[z(1-z)]^{-1/8}/\sqrt2$
is unchanged by $z\mapsto1-z$, proving the Hadamard
fusing relation. It mixes both blocks: the continued
identity block is $(\mathcal F_1+\mathcal F_\epsilon)/\sqrt2$.
Its absolute square therefore cannot equal
$|\mathcal F_1|^2$ for all $z$.



\paragraph{Exercise~\ref{ex:sewing-pairing}.}
\phantomsection
\label{sol:ex:sewing-pairing}
Write $D=\left(\begin{smallmatrix}a&b\\\bar b&d\end{smallmatrix}\right)$.
The relative phase of the two monodromy eigenvalues is $-1$.
With the convention $G=\mathcal F^{\mathsf t}D\overline{\mathcal F}$,
the condition $M^{\mathsf t}D\overline M=D$ forces $b=0$.
For $D=\operatorname{diag}(a,d)$, the condition
$F^{\mathsf t}D\overline F=D$ with the real Hadamard $F$ forces $a=d$.
Thus $D=aI$. Positivity gives $a\ge0$, and a nontrivial
theory with unit vacuum coefficient has $a=1$.
Because $\mathcal F_\epsilon\sim\frac12z^{3/8}$,
the energy-channel bulk coefficient is $1/4$, giving
$|C_{\sigma\sigma\epsilon}|=1/2$.



\paragraph{Exercise~\ref{ex:sewing-object}.}
\phantomsection
\label{sol:ex:sewing-object}
An equivalence of full local theories preserves all
correlators and their sewing structures. An algebra-object
description specifies the chiral multiplication, unit, and
compatibilities that enter the rational full-theory construction.
Forgetting multiplication leaves an underlying module, which
still remembers how the chiral algebra acts. Taking its
character retains only energy multiplicities. Thus OPE
products are lost on passing to the module, and most
representation matrices are lost on passing to the character.
An equality at a weaker level does not reconstruct the stronger
one. In the rational construction, Morita equivalence of
appropriate algebra objects can describe equivalent full
theories even when the algebra objects are not isomorphic.



\paragraph{Exercise~\ref{ex:sewing-coset}.}
\phantomsection
\label{sol:ex:sewing-coset}
If infinitely many orthogonal states have energy $E\le E_*$,
each contributes at least $e^{-\beta E_*}>0$ to the trace.
The partial sums are unbounded for every $\beta>0$.
A common vacuum-energy shift only multiplies this divergence
by a finite positive constant. Positivity rules out cancellations;
an alternating formal series is not a positive thermal trace.



\paragraph{Exercise~\ref{ex:sewing-descent}.}
\phantomsection
\label{sol:ex:sewing-descent}
One needs the chosen reduced Hilbert space and its left-right
pairing, bounded or otherwise controlled line operators on
that space, junction maps with the specified associator,
compatibility with the local OPE, and sewing consistency.
For a nontrivial branching extension, one also needs coherence
with the complementary-sector multiplicities and the algebra
object defining the reduction. A torus trace can detect
spectra and certain line eigenvalue sums. It does not record
the junction associator or all the local and higher-surface
sewing identities.


